%& -shell-escape
% !TeX root = thesis-abstract.tex
% !TeX encoding = UTF-8
% !TeX program = XeLaTeX

\documentclass{.local/lib/classes/ndvr-super-article-standalone}

\begin{document}
 
    При дискретизации в алгоритмах, основанных на $\SAX$,  
    используется равномерная сетка времени.
    Точки смены съёмок в видео расположены неравномерно.
    Кроме того, как было сказано ранее,
    расстояние между соседними съёмками измеряется 
    в относительных единицах.

    В работе предложена модифицированная 
    кусочно-агрегатная аппроксимация. 
    Пусть дан исходный временной ряд $C = (c_{i})^{n}_{i=1}$.
    Ряд $\overline{C}$ строится по следующей формуле:
    \[
        \overline{c}_{i} 
            = \dfrac{w}{n} \cdot
                \sum
                    \limits%
                        ^{\frac{n}{w} \cdot i}%
                        _{j = \frac{n}{w} \cdot (i-1) + 1}%
                c_{j} \cdot \rho_{j}
    \] 
    $\rho_{j}$ — вес съёмки. 
    В качестве веса предлагается 
    использовать вектор относительных длин съёмок.
    $\overline{c}_{i}$ тоже становится вектором.
    Это немного усложняет построение индексного дерева,
    но принцип построения индекса остается неизменным. 

    В алгоритме $\iSAX{2.0}$ величины 
    $n$ и $w$ считаются известными заранее.

    В условиях реального времени такие допущения невозможны.
    Если рассматривать поток видео целиком, $n \to \infty$.  
    Если рассматривать скользящие окна по~потоку,
    $n$ подвержено изменениям.
    Величина $w$ тоже может быть изменена в процессе вычислений.

    В этом случае удобно задать коэффициент дискретизации
    как $\Delta = \frac{n}{w}$.
    \[
        \overline{c}_{i} 
            = \dfrac{1}{\Delta} \cdot
                \sum
                    \limits%
                        ^{\Delta \cdot i}%
                        _{j = \Delta \cdot (i-1) + 1}%
                c_{j} \cdot \rho_{j}
    \] 

    Условия реального времени усложняют 
    построение индекса и поиск по индексу.
    Количество ветвей каждого узла дерева 
    зависит от количества проиндексированных съёмок, 
    что сказывается на объёме индекса и на скорости поиска.
    Но, благодаря настойке параметра $\Delta$,
    есть возможность соблюсти баланс 
    между точностью поиска и объёмом индекса.

\end{document}
